%% Estrategia de aseguramiento de calidad y plan de pruebas
%% Compilar: tectonic -X compile estrategia-qa.tex
\documentclass[11pt,a4paper]{article}
\usepackage{../latex/legasoft}

\lgtitulo{Estrategia de calidad y plan de pruebas}
\lgsubtitulo{Aseguramiento de calidad con ISO/IEC 25010 como modelo de referencia}
\lgtituloCorto{Estrategia de calidad}
\lgcodigo{LGS-CAL-001}
\lgversion{0.1}
\lgestado{Plantilla}
\lgfecha{\campo{fecha}}
\lgautor{\lgQA\ (\lgQARol)}

% --- Ficha reutilizable para un caso de prueba -----------------------------
% Uso: \caso{CP-001}{Título}{RF-001}{Precondiciones}{Pasos}{Resultado esperado}
\newcommand{\caso}[6]{%
\par\vspace{4pt}
\begin{tabularx}{\linewidth}{@{}L{3.4cm} Y@{}}
\toprule
\rowcolor{lgClaro}
\multicolumn{2}{@{}l@{}}{\sffamily\bfseries\color{lgPrimario} #1 \quad #2} \\
\midrule
\sffamily\bfseries\small Requisito cubierto & #3 \\
\sffamily\bfseries\small Precondiciones     & #4 \\
\sffamily\bfseries\small Pasos              & #5 \\
\sffamily\bfseries\small Resultado esperado & #6 \\
\bottomrule
\end{tabularx}
\par\vspace{8pt}}

\begin{document}
\lgportada

\lgcontrolversiones{
0.1 & \campo{dd-mm-aaaa} & \lgQA & Plantilla inicial de estrategia de calidad y plan de pruebas. \\
}

\begin{porcompletar}[Cómo usar esta plantilla]
Este documento reúne la estrategia de calidad y el plan de pruebas, ambos
entregables del \lgQARol\ según \texttt{LGS-INST-001}, sección 10.3. Complete los
campos en ámbar cuando existan requisitos aprobados en \texttt{LGS-REQ-001}: sin
criterios de aceptación no es posible diseñar casos de prueba verificables.
\end{porcompletar}

\tableofcontents
\clearpage

% =============================================================================
\section{Propósito y alcance}

Este documento define cómo Legasoft asegura y comprueba la calidad del producto
\campo{nombre del sistema}. Establece el modelo de calidad adoptado, los niveles
de prueba, los criterios de entrada y salida, la gestión de defectos y la
evidencia que debe conservarse antes de cada entrega.

La función de calidad se ejerce de manera preventiva desde la definición de
requisitos y no solamente al final del desarrollo.

\subsection{Documentos relacionados}

\begin{itemize}
  \item \texttt{LGS-REQ-001} — Especificación de requisitos de software.
  \item \texttt{LGS-ARQ-001} — Documento de arquitectura del sistema.
  \item ISO/IEC 25010 — Modelo de calidad del producto de software.
\end{itemize}

% =============================================================================
\section{Modelo de calidad adoptado}

Legasoft utiliza ISO/IEC 25010 como referencia para definir y evaluar los
atributos del producto. No todas las características aplican con la misma
intensidad a todos los proyectos: la columna de prioridad se ajusta según el
alcance y los riesgos identificados.

\begin{center}
\begin{tabularx}{\linewidth}{@{}L{4.0cm} C{2.0cm} Y@{}}
\toprule
\sffamily\bfseries\color{lgPrimario}Característica &
\sffamily\bfseries\color{lgPrimario}Prioridad &
\sffamily\bfseries\color{lgPrimario}Cómo se evalúa en este proyecto \\
\midrule
Adecuación funcional    & \campo{A/M/B} & Verificación de requisitos y criterios de aceptación. \\
Eficiencia de desempeño & \campo{A/M/B} & \campo{mediciones de tiempo de respuesta o carga.} \\
Compatibilidad          & \campo{A/M/B} & \campo{navegadores, dispositivos o sistemas soportados.} \\
Usabilidad              & \campo{A/M/B} & Revisión de flujos, consistencia y accesibilidad de la interfaz. \\
Fiabilidad              & \campo{A/M/B} & \campo{comportamiento ante fallos y recuperación.} \\
Seguridad               & \campo{A/M/B} & Revisión de control de acceso, datos sensibles y registro de auditoría. \\
Mantenibilidad          & \campo{A/M/B} & Revisiones de código, modularidad y documentación técnica. \\
Portabilidad            & \campo{A/M/B} & \campo{requisitos de despliegue e instalación.} \\
\bottomrule
\end{tabularx}
\captionof{table}{Características de ISO/IEC 25010 aplicadas al proyecto.}
\end{center}

% =============================================================================
\section{Niveles y tipos de prueba}

Estos niveles se reflejan en la estructura de carpetas del repositorio:
\ruta{tests/unit}, \ruta{tests/integration} y \ruta{tests/acceptance}.

\begin{center}
\begin{tabularx}{\linewidth}{@{}L{3.2cm} L{3.4cm} Y L{3.0cm}@{}}
\toprule
\sffamily\bfseries\color{lgPrimario}Nivel &
\sffamily\bfseries\color{lgPrimario}Qué verifica &
\sffamily\bfseries\color{lgPrimario}Responsable y momento &
\sffamily\bfseries\color{lgPrimario}Ubicación \\
\midrule
Unitarias    & Componentes o funciones aisladas.
             & Quien desarrolla el componente, durante la implementación.
             & \ruta{tests/unit} \\
Integración  & Interacción entre componentes, servicios y contratos de API.
             & \lgQARol\ con el \lgArquitectoRol.
             & \ruta{tests/integration} \\
Aceptación   & Cumplimiento de los criterios acordados con el cliente.
             & \lgQARol, antes de cada entrega.
             & \ruta{tests/acceptance} \\
\bottomrule
\end{tabularx}
\captionof{table}{Niveles de prueba y su correspondencia con el repositorio.}
\end{center}

\subsection{Pruebas complementarias}

\begin{itemize}
  \item \textbf{Usabilidad y accesibilidad:} revisión de los flujos de
        interacción y de la adaptación de la interfaz a distintos dispositivos.
  \item \textbf{Regresión:} reejecución del conjunto acordado tras cada cambio
        que afecte a funcionalidad ya validada.
  \item \textbf{Seguridad:} verificación de control de acceso, manejo de datos
        sensibles y trazas de auditoría, dentro del alcance del proyecto.
\end{itemize}

\subsection{Automatización}

La automatización se aplica cuando resulte adecuada y sostenible para el
proyecto. Se prioriza automatizar las pruebas de regresión y los flujos críticos
de negocio; las pruebas exploratorias y de usabilidad permanecen manuales.

\begin{porcompletar}
\campo{Herramientas de prueba seleccionadas. La decisión se registrará como ADR
en docs/arquitectura/ cuando se definan las tecnologías del proyecto.}
\end{porcompletar}

% =============================================================================
\section{Criterios de entrada y salida}

\begin{center}
\begin{tabularx}{\linewidth}{@{}L{3.4cm} Y@{}}
\toprule
\sffamily\bfseries\color{lgPrimario}Criterio &
\sffamily\bfseries\color{lgPrimario}Condición \\
\midrule
Entrada a pruebas & Requisitos con criterios de aceptación aprobados; funcionalidad integrada y desplegada en el entorno de prueba; datos de prueba disponibles. \\
Suspensión        & Un defecto bloqueante impide ejecutar el resto del conjunto, o el entorno de prueba no está operativo. \\
Reanudación       & El defecto bloqueante está corregido y verificado, y el entorno vuelve a estar disponible. \\
Salida de pruebas & Todos los casos de prioridad alta ejecutados y aprobados; sin defectos abiertos de severidad crítica o alta; evidencia registrada e informe de validación emitido. \\
\bottomrule
\end{tabularx}
\captionof{table}{Criterios que gobiernan el ciclo de pruebas.}
\end{center}

% =============================================================================
\section{Diseño de casos de prueba}

Cada caso de prueba se vincula al requisito que verifica, de modo que la matriz
de trazabilidad de \texttt{LGS-REQ-001} pueda completarse sin ambigüedad.

\caso{CP-001}{\campo{título del caso}}
  {\campo{RF-001}}
  {\campo{Estado del sistema y datos necesarios antes de ejecutar.}}
  {\campo{1. \ldots\ 2. \ldots\ 3. \ldots}}
  {\campo{Comportamiento observable que confirma el cumplimiento.}}

\caso{CP-002}{\campo{título del caso}}
  {\campo{RF-002}}
  {\campo{Precondiciones.}}
  {\campo{Pasos.}}
  {\campo{Resultado esperado.}}

% =============================================================================
\section{Gestión de defectos}

Los defectos se registran con información suficiente para su reproducción y se
siguen hasta su resolución. La comunicación es objetiva: describe el
comportamiento observado frente al esperado, sin atribuir responsabilidad.

\subsection{Clasificación por severidad}

\begin{center}
\begin{tabularx}{\linewidth}{@{}L{2.6cm} Y@{}}
\toprule
\sffamily\bfseries\color{lgPrimario}Severidad &
\sffamily\bfseries\color{lgPrimario}Definición \\
\midrule
Crítica & Impide el uso del sistema o compromete la integridad o la seguridad de los datos. Sin solución alternativa. \\
Alta    & Una función principal no cumple su requisito; existe solución alternativa costosa. \\
Media   & Comportamiento incorrecto en una función secundaria, con solución alternativa razonable. \\
Baja    & Defecto cosmético o de redacción que no afecta la operación. \\
\bottomrule
\end{tabularx}
\captionof{table}{Escala de severidad de defectos.}
\end{center}

\subsection{Información mínima de un reporte}

\begin{center}
\begin{tabularx}{\linewidth}{@{}L{3.4cm} Y@{}}
\toprule
\sffamily\bfseries\small\color{lgPrimario}Identificador   & DEF-\campo{NNN} \\
\sffamily\bfseries\small\color{lgPrimario}Caso de prueba  & \campo{CP-NNN} \\
\sffamily\bfseries\small\color{lgPrimario}Severidad       & \campo{crítica / alta / media / baja} \\
\sffamily\bfseries\small\color{lgPrimario}Entorno         & \campo{versión, navegador o dispositivo} \\
\sffamily\bfseries\small\color{lgPrimario}Pasos para reproducir & \campo{secuencia mínima que provoca el defecto} \\
\sffamily\bfseries\small\color{lgPrimario}Resultado obtenido    & \campo{comportamiento observado} \\
\sffamily\bfseries\small\color{lgPrimario}Resultado esperado    & \campo{comportamiento correcto según el requisito} \\
\sffamily\bfseries\small\color{lgPrimario}Evidencia       & \campo{captura, registro o traza} \\
\bottomrule
\end{tabularx}
\captionof{table}{Campos obligatorios de un reporte de defecto.}
\end{center}

% =============================================================================
\section{Métricas de calidad}

Las métricas se emplean para ajustar los procesos y prevenir la repetición de
problemas, en línea con el valor de mejora continua.

\begin{center}
\begin{tabularx}{\linewidth}{@{}L{4.4cm} Y C{2.4cm}@{}}
\toprule
\sffamily\bfseries\color{lgPrimario}Métrica &
\sffamily\bfseries\color{lgPrimario}Cálculo &
\sffamily\bfseries\color{lgPrimario}Meta \\
\midrule
Cobertura de requisitos & Requisitos con al menos un caso de prueba / total de requisitos. & \campo{100\%} \\
Casos ejecutados        & Casos ejecutados / casos planificados.                            & \campo{meta} \\
Tasa de aprobación      & Casos aprobados / casos ejecutados.                               & \campo{meta} \\
Densidad de defectos    & Defectos abiertos / funcionalidades entregadas.                   & \campo{meta} \\
Defectos reabiertos     & Defectos reabiertos / defectos cerrados.                          & \campo{meta} \\
\bottomrule
\end{tabularx}
\captionof{table}{Métricas de seguimiento del proceso de pruebas.}
\end{center}

% =============================================================================
\section{Informe de validación previo a la entrega}

Antes de cada entrega se emite un informe que resume la evidencia y respalda la
decisión de liberar o no el producto.

\begin{itemize}
  \item Alcance de lo probado y versión evaluada.
  \item Resultados por nivel de prueba y por requisito.
  \item Defectos abiertos con su severidad y el riesgo que representan.
  \item Verificación de los criterios de salida.
  \item Recomendación explícita: liberar, liberar con reservas o no liberar.
\end{itemize}

\begin{nota}[Responsabilidad de la recomendación]
La recomendación es técnica y se sustenta en evidencia. La decisión de entrega
corresponde al \lgDirectorRol, que integra el criterio de calidad con el alcance,
las prioridades y los compromisos asumidos con el cliente.
\end{nota}

\end{document}
